% Layout reconstructed from OpenAI's Finite time blowup for Navier--Stokes
% (8 September 2026); the public release does not include a TeX template.
% Reference PDF: https://cdn.openai.com/pdf/32d9f210-8b73-45e0-91bc-82a30aef8a9a/navier-stokes.pdf
\documentclass[11pt,letterpaper,reqno]{amsart}
\usepackage[margin=30mm]{geometry}
\usepackage[T1]{fontenc}
\usepackage{amsmath,amssymb,mathtools}
\usepackage[sc]{mathpazo}
\usepackage{microtype}
\usepackage{etoolbox}
% Match the compact first-page title placement in the reference PDF.
\makeatletter
\patchcmd{\@maketitle}{\global\topskip42\p@\relax}{\global\topskip0\p@\relax}{}{
  \PackageError{DR-Yaglom}{Cannot adjust title placement}{Check the amsart version.}}
\makeatother
\usepackage{enumitem}
\usepackage[colorlinks=true,linkcolor=blue,citecolor=blue,urlcolor=blue]{hyperref}
\hypersetup{
  pdftitle={Exact critical asymptotics for Derrida--Retaux recursions},
  pdfauthor={Ruiqi Ding; Zehua He; Yutao Liang; Dong Wang; Yushu Zheng}
}
\theoremstyle{plain}
\newtheorem{theorem}{Theorem}[section]
\newtheorem{proposition}[theorem]{Proposition}
\newtheorem{lemma}[theorem]{Lemma}
\newtheorem{corollary}[theorem]{Corollary}
\newcommand{\E}{\mathbb E}\newcommand{\Pp}{\mathbb P}
\newcommand{\Zp}{\mathbb Z_{\ge0}}
\newcommand{\dd}{\,\mathrm d}
\newcommand{\norm}[1]{\left\lVert#1\right\rVert}
\newcommand{\floor}[1]{\lfloor#1\rfloor}
\newcommand{\tail}[2]{\mathcal T_{#1}(#2)}
\DeclareMathOperator{\supp}{supp}

\title[Critical Derrida--Retaux recursions]{Exact critical asymptotics for Derrida--Retaux recursions}
\DeclareRobustCommand{\authormark}[1]{\textsuperscript{\MakeLowercase{#1}}}
\author[R. Ding]{Ruiqi Ding\authormark{1,2,a}}
\author[Z. He]{Zehua He\authormark{1,2,b}}
\author[Y. Liang]{Yutao Liang\authormark{1,2,c}}
% Dong Wang: https://arxiv.org/html/2512.20343v1
\author[D. Wang]{Dong Wang\authormark{2,d}}
\author[Y. Zheng]{Yushu Zheng\authormark{3,e}}
\newcommand{\authoremail}[2]{\textsuperscript{#1}\href{mailto:#2}{\textnormal{#2}}}
\newcommand{\frontaffiliations}{%
  \par\vspace{8pt}
  \begingroup
  \normalfont\footnotesize\centering
  \hypersetup{urlcolor=black}
  \textsuperscript{1}\textit{State Key Laboratory of Mathematical Sciences,
  Academy of Mathematics and Systems Science, Chinese Academy of Sciences}\par
  \authoremail{a}{dingruiqi@amss.ac.cn},
  \authoremail{b}{hezehua@amss.ac.cn},
  \authoremail{c}{liangyutao@amss.ac.cn}\par\vspace{5pt}
  \textsuperscript{2}\textit{School of Mathematical Sciences,
  University of Chinese Academy of Sciences}\par
  \authoremail{d}{wangdong@wangd-math.xyz}\par\vspace{5pt}
  \textsuperscript{3}\textit{School of Mathematics,
  Shanghai University of Finance and Economics}\par
  \authoremail{e}{yszheng666@gmail.com}\par
  \endgroup}
\makeatletter
\apptocmd{\@setauthors}{\frontaffiliations}{}{
  \PackageError{DR-Yaglom}{Cannot place affiliations on the first page}{Check the amsart version.}}
\makeatother
\date{}
\begin{document}
\begin{abstract}
We study the critical Derrida--Retaux recursion with a fixed number $m\ge2$ of independent inputs and a nonnegative integer-valued initial law. Under a finite third exponentially tilted moment, excluding the deterministic binary fixed point, we prove that the survival probability is asymptotic to $4/((m-1)^2n^2)$ and identify the locally uniform scaling limit of the positive tilted density, including its boundary value. We also obtain exact asymptotics for the tilted generating function, its one-step decrement and each fixed positive atom, together with convergence of the conditional law to a geometric distribution with total variation error $O(\log n/n)$. The proof combines a cubic moment identity, a truncated convolution estimate and inverse uniqueness for strings of entrance type. A quantitative third-moment tail bound supplies the entrance normalization without a fourth-moment assumption. No fixed smaller polynomial tilted-moment order suffices for these conclusions over all critical initial laws.
\end{abstract}
\maketitle

\section{Introduction}
Recursive models on trees provide a setting in which a phase transition can be studied through the evolution of a single probability law. At criticality, identifying the order of decay of this law and determining its exact asymptotics are distinct questions. In max-type recursions, a small mass near the absorbing state may determine the survival probability even when an exponential tilt places most of the relevant mass on a growing spatial scale.

The Derrida--Retaux model combines independent masses at the children of a vertex, subtracts one unit and replaces negative values by zero. Its binary generating-function iteration was studied by Collet, Eckmann, Glaser and Martin~\cite{CEGM}, and Derrida and Retaux~\cite{DR} later used it as a hierarchical model for the depinning transition. We consider a deterministic number of children and integer-valued masses. The critical condition fixes the first moment of the exponentially tilted law, but does not by itself determine how the positive tilted mass is distributed between the boundary and distances comparable with the generation.

For this recursion, Chen, Derrida, Hu, Lifshits and Shi~\cite{CDHLS} established critical estimates for the generating function and its accumulated product, and formulated more precise asymptotic questions. Under a finite third tilted moment, their estimates control the generating-function excess at order $n^{-1}$. Chen, Hu and Shi~\cite{CHS} subsequently proved that the survival probability and the mean are $n^{-2+o(1)}$, assuming an exponential moment beyond the critical tilting parameter. These results distinguish the decay exponent from the exact boundary constant. Chen and Shi~\cite{CS} showed that heavier tilted tails give a different product exponent, so the moment assumption is relevant to the critical asymptotics. To obtain exact constants from a scaling limit, one must control the density at the boundary as well as at positive macroscopic distances.

We prove a locally uniform limit for the normalized positive tilted density on compact subsets of the closed half-line. Its boundary value yields the survival constant and the asymptotics of every fixed positive atom. The argument applies to every fixed number $m\ge2$ of independent copies in the recursion and uses only a finite third tilted moment. The two estimates that connect the discrete recursion to its limiting profile are a spatial modulus obtained by truncating the convolution evolution and a quantitative third-moment tail bound. They permit the limiting equation to be identified by the inverse uniqueness theorem for strings of entrance type~\cite{Kotani}. We also obtain a quantitative geometric approximation to the conditional law and show that no fixed lower polynomial tilted-moment assumption implies the full collection of asymptotics.

\subsection{Model and main results}
Fix an integer $m\ge2$. Starting from a nonnegative integer-valued random variable $X_0$, define the laws of $X_n$ recursively by
\begin{equation}\label{eq:model}
 X_{n+1}\overset{\mathrm d}=\bigl(X_n^{(1)}+\cdots+X_n^{(m)}-1\bigr)_+,
\end{equation}
where $x_+=\max\{x,0\}$ and, at each step, the variables $X_n^{(1)},\ldots,X_n^{(m)}$ are independent copies of $X_n$. The criticality and moment assumptions are
\begin{equation}\label{eq:assumptions}
 (m-1)\E[X_0m^{X_0}]=\E[m^{X_0}]<\infty,
 \qquad \E[X_0^3m^{X_0}]<\infty.
\end{equation}
When $m=2$, we exclude $\Pp(X_0=1)=1$. This is a deterministic fixed point of the recursion. No other critical deterministic law exists. For $m\ge3$, nonconstant critical laws supported on $\{0,1\}$ are included.

Write
\begin{equation}\label{eq:notation}
 G_n=\E[m^{X_n}],\quad \varepsilon_n=G_n-1,\quad p_n=\Pp(X_n>0),\quad
 q_{n,k}=\frac{m^k\Pp(X_n=k)}{G_n},
\end{equation}
and define
\begin{equation}\label{eq:rho}
 \rho_n(j)=(m-1)q_{n,j+1},\qquad j\in\Zp.
\end{equation}
Thus $q_n$ is a probability law, while $\rho_n$ records its positive atoms after a shift by one. The factor $m-1$ makes the limiting profile independent of $m$.

\begin{theorem}\label{thm:profile}
Assume \eqref{eq:assumptions} and exclude the deterministic binary law specified above. For every $R\in[0,\infty)$,
\begin{equation}\label{eq:mainprofile}
 \sup_{0\le x\le R}\left|n^2\rho_n(\floor{nx})-4e^{-2x}\right|\longrightarrow0.
\end{equation}
Consequently,
\begin{align}
 n(G_n-1)&\longrightarrow\frac2{m-1}, &
 n^2p_n&\longrightarrow\frac4{(m-1)^2},\label{eq:main1}\\
 n^2(\varepsilon_n-\varepsilon_{n+1})&\longrightarrow\frac2{m-1}, &
 n^2\E[X_n]&\longrightarrow\frac{4m}{(m-1)^3}.\label{eq:main2}
\end{align}
For every fixed $k\ge1$,
\begin{equation}\label{eq:atoms}
 n^2\Pp(X_n=k)\longrightarrow \frac4{(m-1)m^k},\qquad
 \E[X_n\mid X_n>0]\longrightarrow\frac m{m-1}.
\end{equation}
Moreover,
\begin{equation}\label{eq:tv}
 \sum_{k\ge1}\left|\Pp(X_n=k\mid X_n>0)-(m-1)m^{-k}\right|
 =O\left(\frac{\log(n+2)}{n+1}\right).
\end{equation}
\end{theorem}

The sum in \eqref{eq:tv} is twice the total variation distance. Its rate concerns the conditional distribution; it does not assert a rate for the convergence of $n^2p_n$. All constants below may depend on $m$ and the fixed initial law, and all limits are taken as $n\to\infty$ unless stated otherwise.

The third-moment assumption cannot be replaced by any fixed smaller polynomial tilted-moment order if the conclusions are required for all critical initial laws.

\begin{theorem}\label{thm:sharpness}
For every fixed $m\ge2$ and every $r\in[0,3)$, there exists a nondegenerate critical initial law with $\E[X_0^rm^{X_0}]<\infty$ but $\E[X_0^3m^{X_0}]=\infty$ for which $n(G_n-1)$ does not converge to $2/(m-1)$. Thus no fixed lower polynomial tilted-moment order suffices for the conclusions of Theorem~\ref{thm:profile} over all critical initial laws.
\end{theorem}

Theorem~\ref{thm:sharpness} concerns a class of moment assumptions. It does not claim that a finite third tilted moment is necessary for each individual law, and it does not classify the borderline power-law tail of exponent four. Its proof is given in Section~\ref{sec:sharpness}.

\subsection{Proof outline}
The proof begins with two coarse critical bounds: $G_n-1=O(n^{-1})$ and an upper bound of order $n^2$ for the product of $G_j^{m-1}$. A cubic polynomial adapted to the tilted recursion converts the product estimate into third-moment control. The coefficient of the linear transport term in the density recursion has a summable defect. The recursion for the zero atom then gives a tail bound of order $n^{-1}$ for the sum of the first positive tilted atoms.

To control individual atoms, we follow coefficients along the linear transport. The nonlinear source becomes a finite sum of convolution powers in these coordinates. Its change between successive generations has a cubic tail estimate. Summing at the moving coefficient visits each source coefficient at most once, which yields
\[
 \rho_n(j)\le C\frac{n+1}{(n+j+1)^3}.
\]
The boundary follows by applying a fixed number of recursion steps. This estimate is obtained before any exact survival asymptotic is used.

The quadratic convolution has a different transport speed from the density itself. Comparing their evolutions produces a telescoping identity for spatial differences. At the third-moment threshold, the remainder has a logarithmic sum. We therefore split the evolution at a time $a$: the part transported from that time has amplitude at most $C(a+1)n^{-3}$, while the later part has gradient at most $Cn^{-3}(1+\log(n/(a+1)))$. For a spatial increment $h$, choosing $a=h$ gives a macroscopic modulus of order $r(1+\log(1/r))$. This yields compactness up to the spatial boundary. Summing a one-step logarithmic gradient bound would not give this modulus.

The remaining issue is to identify a subsequential limit. Weighting the tilted law by the same cubic polynomial gives a Markov chain whose expected positive increment is uniformly bounded. Its tail estimate passes to the limiting density $u$ as
\[
 \int_\delta^\infty x^3u(t,x)\dd x\le C\frac{t^3}{\delta}.
\]
Together with the limiting transport--convolution equation, this estimate controls the small-time remainder in the Laplace transform using only three moments. The moment equations define a string of entrance type, and the remainder determines its singular characteristic on the whole negative axis. Kotani's inverse uniqueness theorem identifies the string; fixing its endpoint and the time normalization then gives $u(t,x)=4t^{-2}e^{-2x/t}$. The boundary value yields the survival and atom limits. The spatial gradient bound gives the error in the conditional geometric approximation, and an exact generating-function identity gives the one-step decrement.

For Theorem~\ref{thm:sharpness}, we choose a critical law with a polynomial tilted tail of exponent $\alpha<4$. The product theorem of Chen and Shi gives exponent $\alpha-2$, whereas the asserted generating-function limit would force exponent two. Choosing $\alpha$ above $r+1$ makes the $r$th tilted moment finite and gives the required contradiction.

\subsection{Further related work}
Related work includes the near-critical free-energy asymptotics
of \cite{CDDHLS}, the continuous-time model of \cite{HMP}, and the
geometric-offspring models of \cite{LZ,AHM}. The latter yield exact
asymptotics within special invariant families. Here the number of
children is fixed, and the initial law need not belong to such a family.

The limiting coalescence equation and its exponential critical profile
already appear in \cite{CDDS}; our contribution concerns convergence
from the discrete recursion, including at the boundary, under the third
tilted-moment assumption. Profile identification uses Kotani's inverse
uniqueness theorem~\cite{Kotani}. We follow his terminology
\emph{string of entrance type} and verify its defining condition
\cite[Section~2, condition~(2.1)]{Kotani} before applying the theorem.

\subsection{Organization}
Section~\ref{sec:prelim} records the tilted recursion, coarse estimates and auxiliary moment identities. Section~\ref{sec:mainproof} proves Theorem~\ref{thm:profile} from the smoothing and profile-identification propositions. Sections~\ref{sec:pointwise} and~\ref{sec:smoothing} establish the discrete pointwise and smoothing estimates. Section~\ref{sec:limits} proves the tail bounds and equations for subsequential limits, and Section~\ref{sec:rigidity} identifies those limits. Section~\ref{sec:sharpness} proves Theorem~\ref{thm:sharpness}.

\subsection{AI-assisted proof development and verification status}
\label{subsec:ai-declaration}

Generative AI was used to produce the mathematical proof draft presented
in this manuscript. OpenAI ChatGPT and Codex also assisted with manuscript
preparation and organization. The authors have manually checked the
mathematical arguments in the current draft. Further revision is planned;
Lean~4 formal verification has not yet been completed.

This version is intended to provide a dated record of the current draft
through an initial GitHub release. It should be read as work in progress,
with the statements and proofs subject to correction. Lean~4 formalization
is planned as part of the subsequent checking process. No completed Lean
verification is claimed for this version. Any future claim of formal
verification will specify the results covered, the correspondence with
the manuscript, and any remaining assumptions or unproved dependencies.

\section{Preliminaries}\label{sec:prelim}
Until Section~\ref{sec:sharpness}, we work under the assumptions of Theorem~\ref{thm:profile}. We collect the exact tilted recursion and the moment estimates used in both the pointwise bound and the smoothing argument. The only estimates used without proof in this section are the two coarse critical bounds in \eqref{eq:coarse}.

For a sequence $f$ on $\Zp$, let $Sf(j)=f(j+1)$, let $Tf(0)=0$ and $Tf(j)=f(j-1)$ for $j\ge1$, and let $\Delta=I-S$. Convolutions are on $\Zp$. Put $L_n=n+1$. For $L\ge1$, write
\[
 \norm f_{\infty,3,L}=\sup_{j\ge0}(1+j/L)^3|f(j)|,
 \qquad \norm f_{1,3,L}=\sum_{j\ge0}(1+j/L)^3|f(j)|.
\]

\subsection{Exact tilted recursion and coarse bounds}
Let $\mathcal G_n(s)=\E[s^{X_n}]$, so that $G_n=\mathcal G_n(m)$. Directly from \eqref{eq:model},
\begin{equation}\label{eq:pgf}
 \mathcal G_{n+1}(s)=\frac{\mathcal G_n(s)^m-\mathcal G_n(0)^m}{s}+\mathcal G_n(0)^m.
\end{equation}
Differentiation at $s=m$ shows that
\[
 m(m-1)\mathcal G'_{n+1}(m)-\mathcal G_{n+1}(m)
 =G_n^{m-1}\{m(m-1)\mathcal G'_n(m)-\mathcal G_n(m)\}.
\]
Thus criticality is preserved and
\begin{equation}\label{eq:qmean}
 \sum_{k\ge0}q_{n,k}=1,\qquad \sum_{k\ge0}kq_{n,k}=\frac1{m-1}.
\end{equation}
The third tilted moment remains finite: if $V$ is the sum of the $m$ inputs, then $X_{n+1}^3m^{X_{n+1}}\le V^3m^V$, and expansion of $V^3$, independence, and the initial third moment give finiteness at every finite generation. This justifies every derivative through order three below.

Put
\[
 H_n(z)=\sum_{k\ge0}q_{n,k}z^k,\quad x_n=q_{n,0},\quad
 \theta_n=1-x_n,\quad D_n=1+(m-1)x_n^m.
\]
Then
\begin{equation}\label{eq:tilt}
 H_{n+1}(z)=\frac{H_n(z)^m+(mz-1)x_n^m}{zD_n},\qquad
 G_{n+1}=\frac{G_n^mD_n}{m}.
\end{equation}
Every apparent division by $z$ is removable coefficientwise.

The zero atom is positive. For $m\ge3$, \eqref{eq:qmean} gives $x_0\ge1-1/(m-1)>0$. For $m=2$, criticality gives
\[
 \Pp(X_0=0)=\sum_{k\ge2}(k-1)2^k\Pp(X_0=k).
\]
If this vanished, the law would be $\delta_1$, which is excluded. The zero-atom recursion preserves positivity. The only deterministic critical law is the excluded binary one, so the nonconstant-law hypothesis of \cite{CDHLS} is satisfied.

The upper bounds in \cite[Theorem~3 and Proposition~1]{CDHLS} therefore apply and give
\begin{equation}\label{eq:coarse}
 0\le\varepsilon_n\le\frac C{n+1},\qquad
 \prod_{j=0}^{n-1}G_j^{m-1}\le C(n+1)^2.
\end{equation}
Changing the endpoint of the product in \cite{CDHLS} changes only a bounded factor. Since $m^k-1\ge m-1$ for $k\ge1$,
\begin{equation}\label{eq:theta}
 p_n\le\frac{\varepsilon_n}{m-1},\qquad
 \theta_n=\frac{\varepsilon_n+p_n}{G_n}=O((n+1)^{-1}).
\end{equation}
In particular $x_n\to1$ and $\inf_n x_n>0$.

\subsection{Moment bounds and transport coefficients}
The next identity is the normalized form of the cubic generating-function identity in \cite[eqs.~(18)--(19)]{CDHLS}. We include a short derivation, since its nonnegative polynomial will also be used to construct the weighted chain in Section~\ref{sec:limits}. Set $\chi_m=1-1/(m-1)$ and
\begin{equation}\label{eq:cubic}
 h_m(k)=k(k+1)\left(k-\frac1{m-1}\right)
 =(k)_3+(3+\chi_m)(k)_2+2\chi_m k,
\end{equation}
where $(k)_r=k(k-1)\cdots(k-r+1)$. All coefficients on the right are nonnegative. The binary polynomial is $k^3-k$, and for $m\ge3$ it is positive at every $k\ge1$.

\begin{lemma}\label{lem:cubic}
Let $J_n=\E_{q_n}h_m(K)$. Then $J_n>0$,
\begin{equation}\label{eq:J}
 D_nJ_{n+1}=mJ_n,\qquad
 J_n=J_0\frac{G_0}{G_n}\prod_{j=0}^{n-1}G_j^{m-1}\le C(n+1)^2.
\end{equation}
Consequently,
\begin{equation}\label{eq:lowmoments}
 \E_{q_n}K^3\le C(n+1)^2,\qquad
 \E_{q_n}K^2\le C(n+1).
\end{equation}
\end{lemma}
\begin{proof}
For an integer $V\ge0$,
\[
 h_m((V-1)_+)=(V)_3+\chi_m(V)_2.
\]
If $V=K_1+\cdots+K_m$ and the inputs have law $q_n$, expansion of the falling factorials, followed by \eqref{eq:qmean}, gives
\[
 \E h_m((V-1)_+)=m\E_{q_n}h_m(K).
\]
The extra zero mass in the tilted recursion contributes nothing. This proves the first identity in \eqref{eq:J}; the second follows from $m/D_n=G_n^m/G_{n+1}$. Positivity is immediate for $m\ge3$. In the binary case the zero-atom argument above forces positive mass above one, and the identity propagates positivity.

For $m\ge3$, $h_m(k)\ge k^3$ for integers $k\ge1$. For $m=2$, $k^3\le(4/3)h_2(k)$ when $k\ge2$, and the contribution at one is bounded. Thus the cubic bound follows from \eqref{eq:coarse}. Cauchy--Schwarz gives $\E K^2\le(\E K\,\E K^3)^{1/2}$, proving the second bound.
\end{proof}

The linear part of the positive-density recursion will have the following coefficients:
\begin{equation}\label{eq:c}
 c_n=\frac{mx_n^{m-1}}{D_n},\qquad C_0=1,\qquad C_n=\prod_{j=0}^{n-1}c_j.
\end{equation}
The polynomial $1+(m-1)x^m-mx^{m-1}$ is nonnegative on $[0,1]$, has a double zero at one, and is $O((1-x)^2)$ there. Hence
\begin{equation}\label{eq:C}
 0<c_n\le1,\qquad 1-c_n=O((n+1)^{-2}),\qquad
 C_n\downarrow C_\infty>0.
\end{equation}
All inverses of $x_n$, $c_n$, and $C_n$ used below are therefore uniformly bounded along the fixed orbit.

We will use the zero-atom recursion to sum the boundary terms in the arrival estimate.
\begin{lemma}\label{lem:atomtail}
The first tilted atom has a summable tail satisfying
\begin{equation}\label{eq:atomtail}
 T_n:=\sum_{s\ge n}q_{s,1}\le\frac C{n+1}.
\end{equation}
\end{lemma}
\begin{proof}
The zero coefficient in \eqref{eq:tilt} is $x_{n+1}=c_n(x_n+q_{n,1})$. Rearranging gives the exact identity
\begin{equation}\label{eq:atomidentity}
 q_{n,1}=\theta_n-\theta_{n+1}+(c_n^{-1}-1)x_{n+1}.
\end{equation}
Sum over $n\le s\le r$. The last summand is nonnegative and at most $C(s+1)^{-2}$. Letting $r\to\infty$, using $\theta_r\to0$, proves \eqref{eq:atomtail}.
\end{proof}

For the density \eqref{eq:rho}, \eqref{eq:qmean}, \eqref{eq:theta}, and \eqref{eq:lowmoments} give
\begin{equation}\label{eq:weightedmass}
 \norm{\rho_n}_{1,3,L_n}\le\frac C{L_n},\qquad
 \sum_{j\ge0}(j+1)\rho_n(j)=1.
\end{equation}
The first inequality follows by expanding the cubic weight; no moment beyond the third is used.

\section{Proof of Theorem~\ref{thm:profile}}\label{sec:mainproof}
We first state the two inputs that connect the discrete density to its asymptotic profile. The first gives compactness including the boundary; the second identifies every possible limit. Their proofs occupy Sections~\ref{sec:pointwise}--\ref{sec:rigidity} and use only the assumptions of Theorem~\ref{thm:profile} and the estimates in Section~\ref{sec:prelim}.

\begin{proposition}\label{prop:smoothing}
Under the assumptions of Theorem~\ref{thm:profile},
\begin{align}
 \norm{\rho_n}_{\infty,3,L_n}&\le C L_n^{-2},\label{eq:sup}\\
 \norm{\Delta\rho_n}_{\infty,3,L_n}&\le C\frac{\log(n+2)}{L_n^3}.\label{eq:gradient}
\end{align}
For every integer $1\le h\le L_n$,
\begin{equation}\label{eq:spmod}
 \norm{\rho_n-S^h\rho_n}_{\infty,3,L_n}
 \le \frac{Ch}{L_n^3}\left(1+\log\frac{L_n}{h}\right).
\end{equation}
The unweighted temporal modulus satisfies the same bound:
\begin{equation}\label{eq:timemod}
 \norm{\rho_{n+h}-\rho_n}_{\infty}
 \le \frac{Ch}{L_n^3}\left(1+\log\frac{L_n}{h}\right).
\end{equation}
\end{proposition}

For each integer $N\ge1$, define the rescaled density on the grid by
\begin{equation}\label{eq:rescaled}
 u_N(n/N,j/N)=N^2\rho_n(j),\qquad n,j\ge0,
\end{equation}
and interpolate bilinearly between grid points. The following proposition is proved in Section~\ref{sec:rigidity}, using the moment and equation estimates of Section~\ref{sec:limits}.

\begin{proposition}\label{prop:identification}
Under the assumptions of Theorem~\ref{thm:profile}, every locally uniform subsequential limit of $u_N$ on $(0,\infty)\times[0,\infty)$ is
\begin{equation}\label{eq:identifiedprofile}
 u(t,x)=\frac4{t^2}e^{-2x/t}.
\end{equation}
\end{proposition}

\begin{proof}[Proof of Theorem~\ref{thm:profile}]
On each rectangle $[\eta,T]\times[0,R]$, with $0<\eta<T<\infty$ and $R<\infty$, Proposition~\ref{prop:smoothing} gives uniform boundedness. Equations~\eqref{eq:spmod} and~\eqref{eq:timemod}, first at grid points and then by interpolation, give a modulus controlled by a constant times $r(1+\log(1/r))$ as $r\downarrow0$. Increments comparable with the time index are controlled by boundedness. The interpolation error at the grid scale is $O(N^{-1}\log(N+2))$ and tends to zero. Arzel\`a--Ascoli and a diagonal extraction therefore give local uniform subsequential compactness on the whole of $(0,\infty)\times[0,\infty)$.

Every subsequence of the interpolated densities has a further locally uniform limit, and Proposition~\ref{prop:identification} identifies every such limit as \eqref{eq:identifiedprofile}. Thus the whole family converges. At $t=1$, the floor-index error is at most
$n^2\norm{\Delta\rho_n}_\infty=O(\log(n+2)/(n+1))$, which tends to zero. This proves \eqref{eq:mainprofile}, including the boundary value
\begin{equation}\label{eq:boundary4}
 n^2\rho_n(0)\longrightarrow4.
\end{equation}

The original survival probability is
\begin{equation}\label{eq:survivalread}
 p_n=\frac{G_n}{m-1}\sum_{j\ge0}m^{-j-1}\rho_n(j).
\end{equation}
Telescoping the spatial differences gives
\[
 \left|p_n-\frac{G_n\rho_n(0)}{(m-1)^2}\right|
 \le\frac{G_n\norm{\Delta\rho_n}_\infty}{(m-1)^3}
 =O\left(\frac{\log(n+2)}{(n+1)^3}\right).
\]
Since $G_n\to1$, \eqref{eq:boundary4} proves the second limit in \eqref{eq:main1}. 
The local profile convergence and \eqref{eq:sup} also imply convergence of the total tilted mass. Indeed, for $A\ge1$,
\[
 n\sum_{j>An}\rho_n(j)\le \frac{C}{(1+A)^2}.
\]
Riemann sums on $[0,A]$, followed by $A\to\infty$, therefore give
$n\sum_j\rho_n(j)\to\int_0^\infty4e^{-2x}\dd x=2$. Hence $n\theta_n\to2/(m-1)$. Since $p_n=O(n^{-2})$ and $\varepsilon_n=G_n\theta_n-p_n$, we obtain $n\varepsilon_n\to2/(m-1)$.

For each fixed $k$,
\[
 \Pp(X_n=k)=\frac{G_n}{m-1}m^{-k}\rho_n(k-1).
\]
Equations \eqref{eq:boundary4} and \eqref{eq:gradient} give the fixed-atom limit. In addition $n^2\Pp(X_n=k)\le Cm^{-k}$ by \eqref{eq:sup}. Dominated convergence therefore yields
\[
 n^2\E[X_n]\longrightarrow\frac4{m-1}\sum_{k\ge1}km^{-k}
 =\frac{4m}{(m-1)^3}.
\]
This proves the mean limit in \eqref{eq:main2}. Dividing by the survival asymptotic gives the conditional mean in \eqref{eq:atoms}.

For the quantitative conditional law, let
$a_n=\sum_{k\ge1}m^{-k}\rho_n(k-1)$. Then
\[
 \left|a_n-\frac{\rho_n(0)}{m-1}\right|\le\frac{\norm{\Delta\rho_n}_\infty}{(m-1)^2},
\]
and
\[
 \left|\rho_n(k-1)-(m-1)a_n\right|
 \le \norm{\Delta\rho_n}_\infty\left(k-1+\frac1{m-1}\right).
\]
Divide by $a_n\asymp n^{-2}$, multiply by $m^{-k}$, and sum over $k$. This proves \eqref{eq:tv} from \eqref{eq:gradient}.

Finally, evaluation of \eqref{eq:pgf} at $m$ gives the exact scalar identity
\[
 m(\varepsilon_n-\varepsilon_{n+1})
 =(m-1)\{1-(1-p_n)^m\}-\{(1+\varepsilon_n)^m-1-m\varepsilon_n\}.
\]
Substitution of the already proved limits yields
\[
 n^2(\varepsilon_n-\varepsilon_{n+1})\longrightarrow
 (m-1)\frac4{(m-1)^2}-\frac{m-1}{2}\frac4{(m-1)^2}
 =\frac2{m-1}.
\]
\end{proof}

\section{Proof of Proposition~\ref{prop:pointwise}}\label{sec:pointwise}
The first step towards Proposition~\ref{prop:smoothing} is a cubic-weighted bound on each atom. We prove it by following the coefficients along the linear transport and estimating the accumulated nonlinear source.

\begin{proposition}\label{prop:pointwise}
For every $n,j\ge0$,
\begin{equation}\label{eq:pointwise}
 \rho_n(j)\le C\frac{n+1}{(n+j+1)^3}.
\end{equation}
This is equivalent to \eqref{eq:sup}.
\end{proposition}


\subsection{Recursion in arrival coordinates}
Let
\[
 Q_s(z)=H_s(z)-x_s,\qquad
 \phi_{s,N}=\frac{q_{s,N-s+1}}{C_s}\quad(N\ge s),\qquad
 F_s(z)=\sum_{N\ge s}\phi_{s,N}z^N.
\]
The index $N=s+k-1$ is preserved when the linear transport moves an atom from $k$ to $k-1$. Dividing by $C_s$ also removes its multiplicative coefficient. Put
\[
 b_s=\phi_{s,s}=\frac{q_{s,1}}{C_s},\qquad
 \kappa_s=\frac{C_s}{x_s},\qquad
 \alpha_{r,s}=\frac1m\binom mr\kappa_s^{r-1},
\]
and define
\begin{equation}\label{eq:arrivalA}
 A_s(z)=\sum_{r=2}^m\alpha_{r,s}z^{(r-2)(1-s)}F_s(z)^r.
\end{equation}
Despite its Laurent notation, each summand has minimum degree $2s+r-2$, hence is an ordinary power series with nonnegative coefficients.

\begin{lemma}\label{lem:arrival}
The exact recursions are
\begin{align}
 F_{s+1}(z)&=F_s(z)-b_sz^s+z^{1-s}A_s(z),\label{eq:arrivalrenewal}\\
 \kappa_{s+1}&=\frac{\kappa_s}{1+\kappa_sb_s}.\label{eq:kappa}
\end{align}
In particular, for $N\ge s$,
\begin{equation}\label{eq:arrivalDuhamel}
 \phi_{s,N}=\phi_{0,N}+\sum_{j=0}^{s-1}[z^{N+j-1}]A_j(z).
\end{equation}
\end{lemma}
\begin{proof}
For $k\ge1$, coefficient extraction from \eqref{eq:tilt} gives
\[
 q_{s+1,k}=c_sq_{s,k+1}+
 \frac1{D_s}[z^{k+1}]\sum_{r=2}^m\binom mr x_s^{m-r}Q_s(z)^r.
\]
Since $Q_s(z)=C_sz^{1-s}F_s(z)$ and $D_sC_{s+1}=mx_s^{m-1}C_s$, the normalized order-$r$ source is exactly the corresponding term in \eqref{eq:arrivalrenewal}. Removing the coefficient at degree $s$ gives the boundary subtraction. Finally
\[
 \frac{C_{s+1}}{x_{s+1}}=\frac{C_s}{x_s+q_{s,1}}
 =\frac{\kappa_s}{1+\kappa_sb_s}.
\]
Iteration at coefficient $N$ proves \eqref{eq:arrivalDuhamel}.
\end{proof}

For $1\le r\le m^2$ the moment bounds imply
\begin{equation}\label{eq:Fmoment}
 F_s(1)\le\frac C{s+1},\qquad
 \sum_{k\ge0}(k+r)^3[z^k]F_s(z)^r\le C_r(s+1)^{3-r}.
\end{equation}
Indeed,
\[
 \sum_{N\ge s}(N+1)^3\phi_{s,N}
 =C_s^{-1}\sum_{k\ge1}(s+k)^3q_{s,k}\le C(s+1)^2,
\]
and $(N_1+\cdots+N_r+r)^3\le r^2\sum_i(N_i+1)^3$. Thus, for $K\ge1$,
\begin{equation}\label{eq:Ftail}
 \sum_{k\ge K}[z^k]F_s(z)^r\le C_r(s+1)^{3-r}K^{-3}.
\end{equation}

\subsection{The source difference estimate}
For a Laurent series $B$, write $\tail K B=\sum_{k\ge K}|[z^k]B|$ whenever the sum is finite. Let
\[
 K_m=2m(m-1)+1.
\]
\begin{lemma}\label{lem:sourcedifference}
For integers $N>s+K_m$,
\begin{equation}\label{eq:sourcedifference}
 \tail{N+s}{A_{s+1}-A_s}
 \le \frac C{N^3}\{1+b_s(s+1)^2\}.
\end{equation}
\end{lemma}
\begin{proof}
Write $F=F_s$, $b=b_s$,
\[
 B=bz^s,\qquad P_r=\alpha_{r,s}z^{(r-1)(1-s)}F^r,
 \qquad G=-B+\sum_{r=2}^mP_r,
\]
so that $F_{s+1}=F+G$. At each outer degree $2\le R\le m$, split the difference into coefficient, shift, and function changes:
\begin{align}\label{eq:threechanges}
 A_{s+1}-A_s=\sum_{R=2}^m\bigl[&
 (\alpha_{R,s+1}-\alpha_{R,s})z^{(R-2)(1-s)}F^R\notag\\
 &+\alpha_{R,s+1}\{z^{-(R-2)s}-z^{(R-2)(1-s)}\}F^R\notag\\
 &+\alpha_{R,s+1}z^{-(R-2)s}\{(F+G)^R-F^R\}\bigr].
\end{align}
All coefficients $\alpha_{R,s}$ are bounded, and \eqref{eq:kappa} implies
$|\alpha_{R,s+1}-\alpha_{R,s}|\le Cb$.

Consider one monomial in the function change. If it selects $n_B$ factors $-B$, $n_F$ factors $F$, and $n_r$ factors $P_r$, put
\[
 e=\sum_{r=2}^m(r-1)n_r,\qquad
 q=n_F+\sum_{r=2}^m rn_r=R-n_B+e.
\]
Its coefficient is bounded by a constant times $b^{n_B}$, its convolution factor is $F^q$, and its Laurent exponent is
\begin{equation}\label{eq:tag}
 \sigma=e+(2-q)s,\qquad 0\le e\le m(m-1),\qquad 0\le q\le m^2.
\end{equation}
Coefficient changes have the same form with $q=R$, $e=R-2$, and one extra factor $b$. The two shift terms have $q=R$ and $e=0$ or $R-2$. The shift change vanishes for $R=2$.

If $q\ge1$, the tail threshold for $F^q$ is
\[
 N+s-\sigma=N+(q-1)s-e\ge N-m(m-1)>N/2.
\]
By \eqref{eq:Ftail}, a monomial with $w$ factors $b$ contributes at most
\begin{equation}\label{eq:tagbound}
 Cb^w(s+1)^{3-q}N^{-3}.
\end{equation}
If $w=0$, a function change contains at least one source factor, since the all-$F$ term was deleted. Thus $q\ge3$. A nonzero shift change also has $q\ge3$, so \eqref{eq:tagbound} is at most $CN^{-3}$. If $w\ge1$, then $b\le C/(s+1)$ and $q+w\ge2$, whence
\[
 b^w(s+1)^{3-q}\le Cb(s+1)^{4-q-w}\le Cb(s+1)^2.
\]
Finally $q=0$ is possible only for the all-$(-B)$ term. Its degree is exactly $2s$, below $N+s$, so it contributes nothing. These cases exhaust the finite expansion \eqref{eq:threechanges}. Their number depends only on $m$, proving the lemma. 
\end{proof}

The source-difference estimate can now be summed along a fixed arrival index.
\begin{proposition}\label{prop:far}
There is a bounded function $\omega(N)\to0$ such that, for $N>s+K_m$,
\begin{equation}\label{eq:far}
 \phi_{s,N}\le\frac{\omega(N)+C(s+1)}{N^3}.
\end{equation}
\end{proposition}
\begin{proof}
Define
\[
 \omega(N)=N^3\left\{\phi_{0,N}+\sum_{k\ge N-1}[z^k]A_0(z)\right\}.
\]
The initial third moment and finite convolution expansion of $A_0$ imply $\omega(N)\to0$ and boundedness. Insert $A_j=A_0+\sum_{v<j}(A_{v+1}-A_v)$ into \eqref{eq:arrivalDuhamel}. The $A_0$ block visits the distinct degrees $N-1,\ldots,N+s-2$ only once. For each fixed $v$, the remaining inner sum is bounded by the absolute tail starting at degree $N+v$. Consequently,
\[
 \phi_{s,N}\le\frac{\omega(N)}{N^3}
 +\frac C{N^3}\left\{s+\sum_{v<s}b_v(v+1)^2\right\}.
\]
Since $b_v\le C_\infty^{-1}q_{v,1}$, exact summation by parts and \eqref{eq:atomtail} give, for $s\ge1$,
\[
 \sum_{v=0}^{s-1}(v+1)^2q_{v,1}
 =T_0+\sum_{v=1}^{s-1}(2v+1)T_v-s^2T_s\le C(s+1).
\]
This proves \eqref{eq:far} for $s\ge1$. For $s=0$, the bound follows directly from the definition of $\omega(N)$.
\end{proof}

\subsection{Proof of Proposition~\ref{prop:pointwise}}
\begin{proof}
If $j>K_m$, apply \eqref{eq:far} with $s=n$, $N=n+j$. Boundedness of $\omega$ and $C_n$ gives \eqref{eq:pointwise}.

There remain only $0\le j\le K_m$. Set $r_0=K_m+2$ and take $n\ge2r_0$. Extracting the positive coefficients in \eqref{eq:tilt} and using \eqref{eq:rho} gives the exact density recursion
\begin{equation}\label{eq:rhor}
 \rho_{s+1}=c_sS\rho_s+\sum_{r=2}^m d_{r,s}T^{r-2}\rho_s^{*r},\qquad
 d_{r,s}=\frac{\binom mr x_s^{m-r}}{D_s(m-1)^{r-1}}.
\end{equation}
Iterate this identity from $s=n-r_0$ to $n$. The homogeneous term at $j$ is at most $\rho_{n-r_0}(j+r_0)=O(n^{-2})$ by \eqref{eq:far}, since $n+j>(n-r_0)+K_m$. Each of the at most $r_0(m-1)$ source terms is at most
\[
 C\norm{\rho_s^{*r}}_\infty\le C\norm{\rho_s}_1^r\le C(n+1)^{-2},
\]
using only \eqref{eq:weightedmass}, $r\ge2$, bounded coefficients, and $c_s\le1$. Thus $\rho_n(j)=O(n^{-2})$ for these finitely many $j$, which is equivalent to \eqref{eq:pointwise} there. For the finitely many generations $n<2r_0$, the finite third moment bounds the cubic-weighted supremum directly. This proves the proposition.
\end{proof}

\section{Proof of Proposition~\ref{prop:smoothing}}\label{sec:smoothing}
Proposition~\ref{prop:pointwise} gives the required pointwise bound. To control spatial differences, we compare the transport of the density with that of its quadratic convolution. We retain a free initial time in the resulting telescope so that the estimate also controls macroscopic increments.
\subsection{Preliminaries}
The weight $(1+j/L)^3$ is submultiplicative. Therefore \eqref{eq:weightedmass} and \eqref{eq:pointwise} imply, for every fixed positive integer $r$,
\begin{equation}\label{eq:convnorms}
 \norm{\rho_s^{*r}}_{1,3,L_s}\le C_rL_s^{-r},\qquad
 \norm{\rho_s^{*r}}_{\infty,3,L_s}\le C_rL_s^{-r-1}.
\end{equation}
The operator $S$ is a contraction in both weighted norms. Each fixed $T^k$ has norm at most $(1+k)^3$. For $0\le s\le n-1$ and $d\ge0$ with $d\ge n-2-s$,
\begin{equation}\label{eq:shiftgain}
 \norm{S^df}_{\infty,3,L_n}\le27\left(\frac{L_s}{L_n}\right)^3
 \norm f_{\infty,3,L_s}\qquad(n\ge2).
\end{equation}
Indeed the ratio of the weights is
$[L_s(n+1+j)/(L_n(s+1+j+d))]^3$, and $s+1+d\ge n-1$.

Put $d_s=d_{2,s}$, $\beta_s=\rho_s(0)$, and
\[
 B_s=\rho_s*\rho_s,\qquad E_s=\sum_{r=3}^m d_{r,s}T^{r-2}\rho_s^{*r}.
\]
The sum is empty when $m=2$. Then
\begin{equation}\label{eq:rhoshort}
 \rho_{s+1}=c_sS\rho_s+d_sB_s+E_s.
\end{equation}
The coefficient limits and variations are
\begin{equation}\label{eq:coefficientbounds}
 d_s\longrightarrow\frac12,\quad |1-c_s|\le CL_s^{-2},\quad
 |d_s-d_{s-1}|\le CL_s^{-2}\quad(s\ge1).
\end{equation}
For the last inequality, $x_{s+1}=c_s(x_s+q_{s,1})$ and Proposition~\ref{prop:pointwise} give $|x_{s+1}-x_s|\le CL_s^{-2}$. The function defining $d_s$ is Lipschitz on the compact range $\inf_sx_s\le x\le1$.

The two coefficientwise convolution identities
\[
 (S\rho)*(S\rho)=S^2(\rho*\rho)-2\rho(0)S^2\rho,
 \qquad (S\rho)*(\rho*\rho)=S(\rho^{*3})-\rho(0)S(\rho*\rho)
\]
give
\begin{equation}\label{eq:Bevolution}
 B_{s+1}=c_s^2S^2B_s+R_s,
\end{equation}
where the complete remainder is
\begin{align}\label{eq:R}
 R_s={}&-2c_s^2\beta_sS^2\rho_s
 +2c_sd_sS(\rho_s^{*3})-2c_sd_s\beta_sSB_s+d_s^2\rho_s^{*4}\notag\\
 &+2(c_sS\rho_s+d_sB_s)*E_s+E_s*E_s.
\end{align}
From \eqref{eq:convnorms} and $\beta_s=O(L_s^{-2})$,
\begin{equation}\label{eq:REbounds}
 \norm{R_s}_{\infty,3,L_s}\le CL_s^{-4},\qquad
 \norm{E_s}_{\infty,3,L_s}\le CL_s^{-4},\qquad
 \norm{E_s}_{1,3,L_s}\le CL_s^{-3}.
\end{equation}
For example, the first term in $R_s$ is $O(L_s^{-2}L_s^{-2})$ and the cubic term is $O(L_s^{-4})$. A term $(S\rho_s)*E_s$ is $O(L_s^{-2}L_s^{-3})$ in the weighted supremum norm; all other terms are no larger.

\subsection{A truncated telescope}
Fix $n\ge4$ and $0\le a\le n-2$. For $a\le s\le n-1$ let
\[
 P_s=\prod_{i=s+1}^{n-1}c_i,\qquad w_s=P_sd_s,\qquad
 K_s=S^{n-1-s}B_s.
\]
Split the exact Duhamel formula as
\begin{align}\label{eq:split}
 \rho_n&=I_{a,n}+V_{a,n},\notag\\
 I_{a,n}&=\frac{C_n}{C_a}S^{n-a}\rho_a,\notag\\
 V_{a,n}&=Q_{a,n}+W_{a,n},\quad
 Q_{a,n}=\sum_{s=a}^{n-1}w_sK_s,\quad
 W_{a,n}=\sum_{s=a}^{n-1}P_sS^{n-1-s}E_s.
\end{align}
Here $I_{a,n}$, $Q_{a,n}$, and $W_{a,n}$ are nonnegative sequences. For $s\le n-2$, \eqref{eq:Bevolution} implies
\[
 K_{s+1}=c_s^2SK_s+S^{n-2-s}R_s,
 \quad \Delta K_s=K_s-c_s^{-2}K_{s+1}+c_s^{-2}S^{n-2-s}R_s.
\]
Consequently the exact truncated identity is
\begin{align}\label{eq:telescope}
 \Delta Q_{a,n}={}&w_aK_a-w_{n-2}c_{n-2}^{-2}B_{n-1}
   +w_{n-1}\Delta B_{n-1}\notag\\
 &+\sum_{s=a+1}^{n-2}(w_s-w_{s-1}c_{s-1}^{-2})K_s\notag\\
 &+\sum_{s=a}^{n-2}w_sc_s^{-2}S^{n-2-s}R_s.
\end{align}
There are no negative shifts in this identity. Since $P_{s-1}=c_sP_s$,
\[
 |w_s-w_{s-1}c_{s-1}^{-2}|
 =P_s|d_s-c_sc_{s-1}^{-2}d_{s-1}|\le CL_s^{-2}.
\]
The two endpoint bounds follow from \eqref{eq:convnorms} and \eqref{eq:shiftgain}. The coefficient-variation sum is bounded by $CL_n^{-3}\sum_{s\ge a+1}L_s^{-2}$. The last sum is bounded by $CL_n^{-3}\sum_{s=a}^{n-2}L_s^{-1}$. Thus
\begin{equation}\label{eq:lateQ}
 \norm{\Delta Q_{a,n}}_{\infty,3,L_n}
 \le \frac C{L_n^3}\left(1+\log\frac{L_n}{a+1}\right).
\end{equation}
Similarly, \eqref{eq:REbounds} and \eqref{eq:shiftgain} give
\[
 \norm{W_{a,n}}_{\infty,3,L_n}
 \le \frac C{L_n^3}\sum_{s=a}^{n-1}\frac1{L_s}.
\]
Since $\norm{\Delta f}_{\infty,3,L}\le2\norm f_{\infty,3,L}$, we conclude that
\begin{equation}\label{eq:lateV}
 \norm{\Delta V_{a,n}}_{\infty,3,L_n}
 \le \frac C{L_n^3}\left(1+\log\frac{L_n}{a+1}\right).
\end{equation}
The logarithm comes from the sum of $L_s^{-1}$ in the remainder. Keeping its dependence on $a$ will be necessary for the spatial modulus.

\subsection{Proof of Proposition~\ref{prop:smoothing}}
By Proposition~\ref{prop:pointwise} and $C_n/C_a\le1$, for $0\le a\le n$,
\[
 I_{a,n}(j)=\frac{C_n}{C_a}\rho_a(j+n-a)
 \le C\frac{a+1}{(n+j+1)^3}.
\]
Hence, in particular for $0\le a\le n/2$,
\begin{equation}\label{eq:earlyI}
 \norm{I_{a,n}}_{\infty,3,L_n}\le C\frac{a+1}{L_n^3}.
\end{equation}
For every integer $h\ge1$, $I-S^h=\sum_{j=0}^{h-1}S^j\Delta$, and $S$ contracts the weighted norm. Equations \eqref{eq:split}, \eqref{eq:lateV}, and \eqref{eq:earlyI} therefore yield
\begin{equation}\label{eq:twoparameter}
 \norm{\rho_n-S^h\rho_n}_{\infty,3,L_n}
 \le \frac C{L_n^3}\left\{a+1+h\left(1+\log\frac{L_n}{a+1}\right)\right\}
\end{equation}
for $0\le a\le n/2$. For $1\le h\le n/2$, choose $a=h$. For $n/2<h\le L_n$, the crude bound $2\norm{\rho_n}_{\infty,3,L_n}\le CL_n^{-2}$ implies the same conclusion. The finitely many small $n$ are absorbed into the constant. This proves \eqref{eq:spmod}; setting $h=1$ proves \eqref{eq:gradient}.

For the temporal bound, iterate \eqref{eq:rhor} from $n$ to $n+h$, with $1\le h\le L_n$. By \eqref{eq:C},
\[
 \left|1-\prod_{s=n}^{n+h-1}c_s\right|\le C\frac h{L_n^2}.
\]
The supremum of each full source in this block is $O(L_n^{-3})$ by \eqref{eq:convnorms}. Consequently,
\[
 \norm{\rho_{n+h}-S^h\rho_n}_\infty\le C\frac h{L_n^3}.
\]
Combine this with \eqref{eq:spmod}. This proves \eqref{eq:timemod} and hence Proposition~\ref{prop:smoothing}.

\section{Moment bounds and equations for subsequential limits}\label{sec:limits}
We now prepare the proof of Proposition~\ref{prop:identification}. The smoothing estimates give local compactness but do not control the third-moment tail. A Markov chain obtained by cubic weighting supplies this missing estimate. We then pass the discrete recursion and its moments to the limit.

\subsection{A cubic-weighted chain and its tail bound}
Recall that $h_m$ is the nonnegative cubic polynomial in \eqref{eq:cubic}, and $J_n=\E_{q_n}h_m(K)$ grows at most quadratically. The next construction realizes the probability laws proportional to $h_mq_n$ on a common Markov chain. It is defined only where $h_m$ is positive.

\begin{lemma}\label{lem:spine}
There is a Markov chain $K_n^\sharp$ with marginals
\[
 q_n^\sharp(k)=\frac{h_m(k)q_{n,k}}{J_n}
\]
such that
\begin{equation}\label{eq:increment}
 \E[(K_{n+1}^\sharp-K_n^\sharp)_+]\le4\qquad(n\ge0).
\end{equation}
In particular, for $R>0$,
\begin{equation}\label{eq:sharptail}
 q_n^\sharp(K>R)\le q_0^\sharp(K>R/2)+\frac{8n}{R}.
\end{equation}
\end{lemma}
\begin{proof}
For $k\in\supp q_n^\sharp$ and $y=(y_1,\ldots,y_{m-1})$, define
\begin{align}\label{eq:allocation}
 A_m(k;y)={}&(k)_3+(k)_2\left(3\sum_i y_i+\chi_m\right)\notag\\
 &+k\left(2\sum_{i<j}y_iy_j+\chi_m\sum_i y_i\right).
\end{align}
All terms are nonnegative. Choose the companions by weighting the product law $q_n^{\otimes(m-1)}$ with the density
\[
 \frac{A_m(k;y)}{h_m(k)}.
\]
This density integrates to one: each companion has mean $1/(m-1)$, so integration of \eqref{eq:allocation} gives $h_m(k)$. Set
\[
 K_{n+1}^\sharp=k+y_1+\cdots+y_{m-1}-1.
\]
The factorial expansion gives, for $V=k_1+\cdots+k_m$,
\[
 \sum_{i=1}^m A_m(k_i;(k_j)_{j\ne i})=h_m((V-1)_+).
\]
Averaging the distinguished carrier contributes $1/m$. Thus the mass at an output $\ell\ge1$ is
$h_m(\ell)\Pp_{q_n}(V-1=\ell)/(mJ_n)=q_{n+1}^\sharp(\ell)$ by \eqref{eq:J}. At any output with $h_m(\ell)=0$, the sum of the nonnegative allocations is zero, so no mass is lost. In particular the binary support stays above one, as required for the transition density to be defined.

For completeness, put $b=\E_{q_n}(K)_2$, $g=\E_{q_n}(K)_3$, $x=x_n$, and $\chi=\chi_m$. Integrating the positive increment with the allocation density gives the exact expression
\begin{equation}\label{eq:incformula}
 \E[(K_{n+1}^\sharp-K_n^\sharp)_+]
 =\frac{x^{m-1}(g+\chi b)+3(m-1)b^2+6\chi b+3\chi^2/(m-1)}
 {g+(3+\chi)b+2\chi/(m-1)}.
\end{equation}
Here is a direct calculation. For independent companions, let $U=\sum_iY_i$ and $V=\sum_{i<j}Y_iY_j$. Then
\[
 \E(U-1)_+=x^{m-1},\quad \E[U(U-1)]=(m-1)b+\chi,
\]
\[
 \E[V(U-1)]=(m-2)b+\chi^2.
\]
Multiplying these identities by the integrated carrier factors $g$, $b$, and $1/(m-1)$ in \eqref{eq:allocation} gives \eqref{eq:incformula}.

Cauchy--Schwarz for the size-biased law gives
$b^2\le(g+b)/(m-1)$. Since $0\le\chi<1$ and $x\le1$, the numerator in \eqref{eq:incformula} is at most
\[
 4g+(3+7\chi)b+\frac{3\chi^2}{m-1}
 \le4\left\{g+(3+\chi)b+\frac{2\chi}{m-1}\right\}.
\]
This proves \eqref{eq:increment}, including $m=2$ where $\chi=0$.

Pathwise, $K_n^\sharp\le K_0^\sharp+\sum_{s<n}(K_{s+1}^\sharp-K_s^\sharp)_+$. A union bound and Markov's inequality yield \eqref{eq:sharptail}. The estimate only uses the tail of $K_0^\sharp$; its first moment need not be finite.
\end{proof}

The weighted-chain estimate converts the cubic moment bound into uniform integrability on the scaling used in \eqref{eq:rescaled}.
\begin{corollary}\label{cor:thirdtail}
For $R\ge2$,
\begin{equation}\label{eq:thirdtail}
 \sum_{k>R}k^3q_{n,k}
 \le C(n+1)^2\left\{q_0^\sharp(K>R/2)+\frac nR\right\}.
\end{equation}
For every $0<\eta<T<\infty$,
\begin{equation}\label{eq:UI}
 \lim_{A\to\infty}\limsup_{N\to\infty}
 \sup_{\eta N\le n\le TN}N^{-2}\sum_{k>AN}k^3q_{n,k}=0.
\end{equation}
\end{corollary}
\begin{proof}
For $k>R\ge2$, $k^3\le C_mh_m(k)$ for every fixed $m\ge2$. Multiply \eqref{eq:sharptail} by $C_mJ_n$ and use \eqref{eq:J}. For \eqref{eq:UI}, divide by $N^2$, take the supremum, first let $N\to\infty$ to remove the tail of the single fixed probability law $q_0^\sharp$, and then let $A\to\infty$. The order of limits is essential.
\end{proof}

\subsection{Proof of Proposition~\ref{prop:limit}}
Let $u$ be a locally uniform subsequential limit of the interpolated densities in \eqref{eq:rescaled}. Such limits exist by Proposition~\ref{prop:smoothing}, as shown at the start of the proof of Theorem~\ref{thm:profile}. We work along the corresponding subsequence throughout this section.

\begin{proposition}\label{prop:limit}
The limit $u$ has moments
\[
 A_r(t)=\int_0^\infty x^ru(t,x)\dd x\qquad(0\le r\le3)
\]
satisfying
\begin{equation}\label{eq:Amoment}
 A_0(t)\le C/t,\qquad A_1(t)=1,\qquad A_2(t)\le Ct,\qquad A_3(t)\le Ct^2.
\end{equation}
The third moment also satisfies the small-time tail estimate
\begin{equation}\label{eq:limitthirdtail}
 \int_{x>\delta}x^3u(t,x)\dd x\le C\frac{t^3}{\delta}\qquad(t,\delta>0).
\end{equation}
The mild equation is
\begin{equation}\label{eq:mild}
 u(t,x)=u(t_0,x+t-t_0)+\frac12\int_{t_0}^t(u(s,\cdot)*u(s,\cdot))(x+t-s)\dd s
\end{equation}
for every $0<t_0<t$ and $x\ge0$.
\end{proposition}
\begin{proof}
For a grid time with $n/N\to t>0$, consider the measure
\[
 \mu_{N,t}=N\sum_{j\ge0}\rho_n(j)\delta_{j/N}.
\]
Its local limit is $u(t,x)\dd x$, by locally uniform convergence and the Riemann-sum normalization. Its $r$th moment is $N^{1-r}\sum_jj^r\rho_n(j)$. The coarse bounds and Corollary~\ref{cor:thirdtail} give uniform integrability through order three on every time annulus. Hence moments pass to the limit:
\begin{equation}\label{eq:momentpass}
 A_r(t)=\lim_{N\to\infty}N^{1-r}\sum_{j\ge0}j^r\rho_n(j),\quad0\le r\le3.
\end{equation}
For $r=0$ the summand has mass weight one. At a non-grid time one may use either adjacent grid time: both approach the same positive time, and their locally uniform limits agree. Equivalently, temporal interpolation gives a convex combination of these discrete measures. Spatial interpolation does not change the limit, by the compact Riemann-sum estimate and the same tail control. The bounds in \eqref{eq:Amoment} follow from \eqref{eq:weightedmass} and \eqref{eq:lowmoments}; the first-moment identity is
\[
 \sum_jj\rho_n(j)=1-\sum_j\rho_n(j)=1-(m-1)\theta_n\longrightarrow1.
\]

For the more precise tail, \eqref{eq:thirdtail} gives, before taking limits,
\[
 N^{-2}\sum_{j>\delta N}j^3\rho_n(j)
 \le C\left(\frac{n+1}{N}\right)^2
 \left\{q_0^\sharp(K>\delta N/2)+\frac n{\delta N}\right\}.
\]
Portmanteau with bounded lower-semicontinuous truncations, followed by monotone convergence, proves \eqref{eq:limitthirdtail}. The initial tail disappears only as $N\to\infty$. For $t$ in a fixed time annulus, the right side is uniformly $O(1/\delta)$.

Now iterate \eqref{eq:rhor} between $\ell=\floor{Nt_0}$ and $n=\floor{Nt}$. The product of unary coefficients tends uniformly to one because $\sum_{s=\ell}^{n-1}(1-c_s)=O(N^{-1})$. The quadratic coefficient tends to $1/2$. On compact space-time sets with positive lower time endpoint,
\[
 N^3(\rho_s*\rho_s)(k)
 =\frac1N\sum_{i=0}^k[N^2\rho_s(i)][N^2\rho_s(k-i)]
\]
converges locally uniformly to the corresponding continuum convolution. For $r\ge3$, \eqref{eq:convnorms} gives a total contribution after time summation and multiplication by $N^2$ bounded by $CN^{2-r}$; each such contribution tends to zero. Fixed right shifts $T^{r-2}$ do not alter that bound. The quadratic space and time Riemann sums therefore give \eqref{eq:mild}. Establishing it first on a countable dense set and using continuity yields every $t_0,t,x$ as stated.
\end{proof}

The pointwise bound also passes to the limit as $u(t,x)\le Ct(t+x)^{-3}$. This bound alone would not control the third-moment tail, so the weighted-chain estimate is needed for the moment passage.

\subsection{Proof of Lemma~\ref{lem:momentODE}}
The mild equation gives the time regularity needed to construct the string of entrance type. We derive the moment equations in weak form, using the continuous boundary trace and the tail estimate from Proposition~\ref{prop:limit}.

\begin{lemma}\label{lem:momentODE}
The functions $A_0,A_2,A_3$ and, for each $p\ge0$,
\[
 U(t,p)=\int_0^\infty e^{-px}u(t,x)\dd x
\]
are locally absolutely continuous. With $\beta(t)=u(t,0)$,
\begin{align}
 A_0'&=-\beta+\tfrac12 A_0^2,&
 A_2'&=A_0A_2-1,& A_3'&=A_0A_3,\label{eq:momentODE}\\
 \partial_tU(t,p)&=pU(t,p)+\tfrac12U(t,p)^2-\beta(t).\label{eq:laplaceODE}
\end{align}
\end{lemma}
\begin{proof}
The mild equation implies $u_t=u_x+\tfrac12u*u$ in distributions. More explicitly, for $\varphi\in C_c^1([0,\infty))$, translating in \eqref{eq:mild} and integrating gives the time-distribution identity
\begin{align*}
 \frac{\mathrm d}{\mathrm dt}\int\varphi(x)u(t,x)\dd x
 ={}&-\varphi(0)u(t,0)-\int\varphi'(x)u(t,x)\dd x\\
 &+\frac12\iint\varphi(x+y)u(t,x)u(t,y)\dd x\dd y.
\end{align*}
Continuity of $u$ up to zero identifies the boundary term. All terms are locally integrable in time by \eqref{eq:Amoment}.

Take smooth cutoffs of $1,x^2,x^3$ and $e^{-px}$. For polynomial degree $r\le3$, choose them with derivative bounded by $C_rx^{r-1}$ away from the mass case, uniformly in the cutoff radius. The transport terms are dominated by moments of order at most two. The source terms are dominated by $C_r(x^r+y^r)$ and are integrable by \eqref{eq:Amoment}. The left sides converge uniformly on time annuli by \eqref{eq:limitthirdtail}. Removing the cutoffs proves absolute continuity and the stated differential identities. In the convolution moments, $A_1=1$ cancels the $A_2$ term in the third equation and gives $A_2'=A_0A_2-1$. Local continuity of the moments follows from continuity of $u$ and the uniform tail bound on time annuli. The right sides of these equations are therefore continuous, so the absolutely continuous functions are in fact continuously differentiable. No differentiability of $u$ in space is required.
\end{proof}

\section{Proof of Proposition~\ref{prop:identification}}\label{sec:rigidity}
Fix a subsequential limit $u$. Recall that its first moment is one, its third moment is $O(t^2)$, and its third-moment tail satisfies \eqref{eq:limitthirdtail}. We use these facts to determine a singular characteristic from the moment and Laplace-transform equations. The construction below verifies the endpoint conditions and the canonical normalization required for inverse uniqueness.

\subsection{Proof of Lemma~\ref{lem:remainder}}
By H\"older and Cauchy--Schwarz,
\begin{equation}\label{eq:A3lower}
 1=A_1^3\le A_0^2A_3,\qquad A_2^2\le A_1A_3=A_3.
\end{equation}
Thus $ct^2\le A_3(t)\le Ct^2$ and $A_2(t)\le Ct$. For $p>0$, set
\[
 Z(t,p)=p-A_0(t)+U(t,p)=\int(e^{-px}-1+px)u(t,x)\dd x.
\]
Then $Z>0$, $Z\le p^2A_2/2=O_p(t)$, and cancellation of the boundary term in \eqref{eq:laplaceODE} gives
\begin{equation}\label{eq:ZODE}
 \partial_tZ=A_0Z+\tfrac12(Z^2-p^2).
\end{equation}

\begin{lemma}\label{lem:remainder}
For each fixed $p>0$,
\begin{equation}\label{eq:taylorthird}
 Z(t,p)=\frac{p^2}{2}A_2(t)-\frac{p^3}{6}A_3(t)+o(A_3(t))\qquad(t\downarrow0).
\end{equation}
In fact the remainder divided by $A_3(t)$ is $O_p(\sqrt t)$ for $0<t\le1$.
\end{lemma}
\begin{proof}
The nonnegative remainder
$R_p(x)=e^{-px}-1+px-p^2x^2/2+p^3x^3/6$ satisfies
\[
 0\le R_p(x)\le \min\{p^4x^4/24,\ p^3x^3/6\}.
\]
Split its integral at a number $\delta>0$. By \eqref{eq:limitthirdtail} and \eqref{eq:A3lower},
\[
 \frac{\int R_p(x)u(t,x)\dd x}{A_3(t)}
 \le C_p\delta+C_p\frac{t}{\delta}.
\]
Taking $\delta=\sqrt t$ proves the assertion. The quantitative tail in \eqref{eq:limitthirdtail} is the input that makes the normalized remainder vanish.
\end{proof}

\subsection{Construction of the string of entrance type}
Fix a normalization constant $\kappa>0$ and put
\begin{equation}\label{eq:Y}
 Y(t)=\frac{\kappa}{A_3(t)}.
\end{equation}
Then $Y'=-A_0Y$, $Y(t)\asymp t^{-2}$, and $(A_2Y)'=-Y$. Since $A_2(t)Y(t)=O(t^{-1})\to0$ as $t\to\infty$,
\[
 A_2(t)Y(t)=\int_t^\infty Y(s)\dd s.
\]
Define
\begin{equation}\label{eq:stringcoordinates}
 \xi(t)=-\int_t^\infty Y(s)\dd s=-A_2(t)Y(t),\qquad
 \mathfrak m(\xi(t))=\int_0^tY(s)^{-1}\dd s.
\end{equation}
The map $t\mapsto\xi(t)$ is increasing from $(0,\infty)$ onto $(-\infty,0)$, and
\begin{equation}\label{eq:stringdensity}
 \xi'(t)=Y(t),\qquad \mathfrak m'(\xi(t))=Y(t)^{-2}.
\end{equation}
Extend $\mathfrak m(\xi)=\infty$ for $\xi\ge0$. The estimates above give
\[
 -\xi(t)\asymp t^{-1},\qquad \mathfrak m(\xi(t))\asymp t^3,
 \qquad \mathfrak m'(\xi)\asymp|\xi|^{-4}\quad(\xi\to-\infty).
\]
Consequently $\mathfrak m(-\infty)=0$, the support is nonempty, the finite right endpoint is zero, and for every $b<0$,
\begin{equation}\label{eq:stringintegrability}
 \int_{-\infty}^b|\xi|\dd\mathfrak m(\xi)<\infty,
 \qquad \int_{-\infty}^b\xi^2\dd\mathfrak m(\xi)<\infty.
\end{equation}
The first condition in \eqref{eq:stringintegrability} is Kotani's entrance condition \cite[condition~(2.1)]{Kotani}; thus $\mathfrak m$ is a \emph{string of entrance type} in his terminology. The second is the stronger condition used to define the singular characteristic below.
Also $\mathfrak m(0-)=\int_0^\infty Y^{-1}\dd t=\infty$. Thus the right-boundary nonuniqueness alternative in Kotani's equation (2.3) does not occur.

We use the following precise published input. For nondecreasing strings of entrance type satisfying the first condition in \eqref{eq:stringintegrability}, equality of canonical spectral measures implies equality up to translation \cite[Theorem 1, arXiv p.~4]{Kotani}. Under the stronger second condition, the singular characteristic is
\begin{equation}\label{eq:characteristic}
 h(\lambda)=\lim_{\xi\to-\infty}\left\{\xi+\varphi_\lambda(\xi)
 \int_\xi^0\varphi_\lambda(v)^{-2}\dd v\right\},\qquad\lambda<0,
\end{equation}
where $\varphi_\lambda$ is the solution normalized by $\varphi_\lambda(-\infty)=1$. Its Herglotz representation on arXiv p.~2 identifies its spectral measure. Therefore equality of these characteristics for every negative parameter implies equality of the spectral measures. Fixing the right endpoint at zero removes the translation freedom.

For $p>0$ put
\begin{equation}\label{eq:Phi}
 \Phi(t,p)=\exp\left\{-\frac12\int_0^tZ(s,p)\dd s\right\},\qquad
 f_p(\xi(t))=\frac2{p^2}Y(t)Z(t,p)\Phi(t,p).
\end{equation}
The bounds on $Z$ imply $0<\Phi\le1$ and $\Phi=1+O_p(t^2)$ near zero. Direct differentiation using \eqref{eq:ZODE} gives
\begin{equation}\label{eq:fODE}
 \frac{\mathrm d}{\mathrm dt}f_p(\xi(t))=-Y(t)\Phi(t,p),\qquad
 f'_p(\xi(t))=-\Phi(t,p),\qquad
 f''_p=\frac{p^2}{4}\mathfrak m'f_p.
\end{equation}
The solution is positive, decreasing, and satisfies $f_p(0-)=0$, since
$f_p(\xi(t))\le Y(t)A_2(t)=O(t^{-1})$ as $t\to\infty$.

Lemma~\ref{lem:remainder}, together with $YA_3=\kappa$, gives
\begin{align}\label{eq:fpentrance}
 f_p(\xi(t))+\xi(t)
 &=YA_2(\Phi-1)-\frac{\kappa p}{3}\Phi
   +\frac2{p^2}Y\Phi\int R_p(x)u(t,x)\dd x\notag\\
 &=-\frac{\kappa p}{3}+O_p(\sqrt t).
\end{align}
Furthermore $f'_p(\xi)\to-1$ as $\xi\to-\infty$.

\subsection{Canonical normalization and inverse uniqueness}
Let $q=p^2/4$ and let $\varphi=\varphi_{-q}$ be the solution normalized at $-\infty$. Its Volterra equation is
\[
 \varphi(\xi)=1+q\int_{-\infty}^\xi(\xi-v)\varphi(v)\mathfrak m'(v)\dd v.
\]
Thus $\varphi\ge1$ and $\varphi'\ge0$. Since
$B(\xi)=\int_{-\infty}^\xi(\xi-v)\mathfrak m'(v)\dd v=O(|\xi|^{-2})$, monotonicity gives
$\varphi(\xi)\le1+qB(\xi)\varphi(\xi)$. For sufficiently negative $\xi$,
\[
 \varphi(\xi)=1+O_q(|\xi|^{-2}),\qquad
 \varphi'(\xi)=O_q(|\xi|^{-3}).
\]
Equations \eqref{eq:fODE} and \eqref{eq:fpentrance} then imply that the constant Wronskian is
\[
 f_p\varphi'-f'_p\varphi=1.
\]
Since $f_p/\varphi\to0$ at $0-$, integration of $(f_p/\varphi)'=-\varphi^{-2}$ proves
\[
 f_p(\xi)=\varphi(\xi)\int_\xi^0\varphi(v)^{-2}\dd v.
\]
This identifies $f_p$ with the canonical solution in \eqref{eq:characteristic}. From \eqref{eq:fpentrance},
\begin{equation}\label{eq:h}
 h(-q)=-\frac{2\kappa}{3}\sqrt q\qquad(q>0).
\end{equation}
This is an identity on the whole negative axis.

Put $d=2\kappa/3$ and consider the explicit string
\begin{equation}\label{eq:modelstring}
 \mathfrak m_*(\xi)=\frac{d^2}{3(-\xi)^3}\quad(\xi<0),\qquad
 \mathfrak m_*(\xi)=\infty\quad(\xi\ge0).
\end{equation}
Its decreasing normalized solution is
\[
 f_{*,p}(\xi)=(-\xi)\exp\left\{\frac{dp}{2\xi}\right\}.
\]
A direct derivative calculation gives $f''_{*,p}=(p^2/4)\mathfrak m_*'f_{*,p}$, $f_{*,p}(0-)=0$, $f'_{*,p}(-\infty)=-1$, and
\[
 \lim_{\xi\to-\infty}\{f_{*,p}(\xi)+\xi\}=-dp/2=-\kappa p/3.
\]
The same Wronskian normalization applies, so its characteristic is exactly \eqref{eq:h}. Both strings meet \eqref{eq:stringintegrability}. Inverse uniqueness and their common finite right endpoint give $\mathfrak m=\mathfrak m_*$.

To recover physical time rather than just a reparametrized string, compare their densities:
\[
 Y(t)^{-2}=d^2\xi(t)^{-4},\qquad \xi'(t)=Y(t).
\]
Thus $\xi'=\xi^2/d$. The condition $\xi(t)\to-\infty$ as $t\downarrow0$ fixes the integration constant, giving
\[
 \xi(t)=-d/t,\qquad Y(t)=d/t^2,\qquad A_0(t)=-Y'(t)/Y(t)=2/t.
\]
Uniqueness of the normalized decreasing solution yields $f_p=f_{*,p}$. From $f'_p(\xi(t))=-\Phi(t,p)$ one obtains
\[
 \Phi(t,p)=(1+pt/2)e^{-pt/2},\qquad
 Z(t,p)=-2\partial_t\log\Phi(t,p)=\frac{tp^2}{tp+2}.
\]
Therefore
\begin{equation}\label{eq:rigidprofile}
 U(t,p)=A_0(t)-p+Z(t,p)=\frac4{t(tp+2)},\qquad
 u(t,x)=\frac4{t^2}e^{-2x/t}.
\end{equation}
Uniqueness of Laplace transforms proves the last equality. The time variable is the original scaling $t=n/N$. This proves Proposition~\ref{prop:identification}.

\section{Proof of Theorem~\ref{thm:sharpness}}\label{sec:sharpness}
We construct a critical law with a polynomial tilted tail and apply the stable-product theorem of Chen and Shi~\cite[Theorem~1.1]{CS}. The resulting product exponent is incompatible with the limit for $n(G_n-1)$ in Theorem~\ref{thm:profile}.

\begin{proof}
Fix $m\ge2$ and $r\in[0,3)$. Choose $\alpha\in(\max\{3,r+1\},4)$ and $c>0$ sufficiently small. Define a tilted probability law by
\[
 q_{0,k}=ck^{-\alpha}\ (k\ge2),\quad
 q_{0,1}=\frac1{m-1}-c\sum_{k\ge2}k^{1-\alpha},
\]
\[
 q_{0,0}=1-\frac1{m-1}+c\sum_{k\ge2}(k-1)k^{-\alpha}.
\]
All masses are nonnegative, $q_{0,0}>0$, the total mass is one, and the mean is $1/(m-1)$. Put
\[
 G_0=\left(\sum_{k\ge0}m^{-k}q_{0,k}\right)^{-1},\qquad
 \Pp(X_0=k)=G_0m^{-k}q_{0,k}.
\]
The original law is exactly critical, nondegenerate, and has positive mass above one. Its tilted moment of positive order $s$ is finite exactly when $s<\alpha-1$; its zeroth tilted moment is finite as well. Hence the stated $r$th moment is finite and the third is infinite.

The stable-product theorem applies because the original tail is exactly $G_0c\,m^{-k}k^{-\alpha}$ for $k\ge2$. It gives
\[
 \prod_{j=0}^{n-1}G_j^{m-1}\asymp n^{\alpha-2}.
\]
If $n(G_n-1)\to2/(m-1)$, then $\sum_j(G_j-1)^2<\infty$ and
\[
 (m-1)\sum_{j<n}\log G_j=2\log n+o(\log n),
\]
contradicting the product exponent $\alpha-2<2$. This proves the theorem.
\end{proof}
% Candidate funding acknowledgment, pending confirmation that these grants
% apply to the present work. Source: the Acknowledgements in
% https://arxiv.org/html/2307.03720v5 (23 October 2025).
% The supplied four-author manuscript contains no funding statement.
% \section*{Funding}
% Dong Wang acknowledges partial support from NSFC grants 12271502 and
% 11871425 and UCAS start-up grant 118900M043.

\begin{thebibliography}{99}
\bibitem{CDHLS} X. Chen, B. Derrida, Y. Hu, M. Lifshits and Z. Shi (2019).
A max-type recursive model: some properties and open questions.
In \emph{Sojourns in Probability Theory and Statistical Physics--III},
Springer Proceedings in Mathematics \& Statistics \textbf{300}, 166--186.
\href{https://arxiv.org/abs/1705.04787v2}{arXiv:1705.04787v2}.

\bibitem{CHS} X. Chen, Y. Hu and Z. Shi (2022).
The sustainability probability for the critical Derrida--Retaux model.
\emph{Probab. Theory Related Fields} \textbf{182}, 641--684.
\href{https://doi.org/10.1007/s00440-021-01091-z}{doi:10.1007/s00440-021-01091-z}.

\bibitem{CS} X. Chen and Z. Shi (2021).
The stable Derrida--Retaux system at criticality.
In \emph{In and Out of Equilibrium 3: Celebrating Vladas Sidoravicius},
Progress in Probability \textbf{77}, 239--264.
\href{https://arxiv.org/abs/2006.06140v1}{arXiv:2006.06140v1}.

\bibitem{CEGM} P. Collet, J.-P. Eckmann, V. Glaser and A. Martin (1984).
Study of the iterations of a mapping associated to a spin glass model.
\emph{Comm. Math. Phys.} \textbf{94}, 353--370.
\href{https://doi.org/10.1007/BF01224830}{doi:10.1007/BF01224830}.

\bibitem{DR} B. Derrida and M. Retaux (2014).
The depinning transition in presence of disorder: a toy model.
\emph{J. Stat. Phys.} \textbf{156}, 268--290.
\href{https://arxiv.org/abs/1401.6919}{arXiv:1401.6919}.

\bibitem{Kotani} S. Kotani (2013).
Krein's strings whose spectral functions are of polynomial growth.
\emph{Kyoto J. Math.} \textbf{53}, 787--814.
\href{https://arxiv.org/abs/1304.6786v1}{arXiv:1304.6786v1}.

\bibitem{CDDHLS} X. Chen, V. Dagard, B. Derrida, Y. Hu, M. Lifshits and Z. Shi (2021).
The Derrida--Retaux conjecture on recursive models.
\emph{Ann. Probab.} \textbf{49}, 637--670.
\href{https://doi.org/10.1214/20-AOP1457}{doi:10.1214/20-AOP1457}.

\bibitem{CDDS} X. Chen, V. Dagard, B. Derrida and Z. Shi (2020).
The critical behaviors and the scaling functions of a coalescence equation.
\emph{J. Phys. A: Math. Theor.} \textbf{53}, 195202.
\href{https://arxiv.org/abs/2001.00853}{arXiv:2001.00853}.

\bibitem{HMP} Y. Hu, B. Mallein and M. Pain (2020).
An exactly solvable continuous-time Derrida--Retaux model.
\emph{Comm. Math. Phys.} \textbf{375}, 605--651.
\href{https://doi.org/10.1007/s00220-019-03465-w}{doi:10.1007/s00220-019-03465-w}.

\bibitem{LZ} Z. Li and R. Zhang (2025).
Asymptotic behavior of the generalized Derrida--Retaux recursive model.
\emph{J. Stat. Phys.} \textbf{192}, article 98.
\href{https://doi.org/10.1007/s10955-025-03480-3}{doi:10.1007/s10955-025-03480-3}.

\bibitem{AHM} G. Alsmeyer, Y. Hu and B. Mallein (2025).
The Derrida--Retaux model on a geometric Galton--Watson tree.
Preprint.
\href{https://arxiv.org/abs/2502.02991}{arXiv:2502.02991}.
\end{thebibliography}
\end{document}
